\documentclass[11pt,a4paper]{article}
\usepackage[margin=1in]{geometry}
\usepackage{booktabs}
\usepackage{hyperref}
\usepackage{xcolor}
\usepackage{parskip}
\usepackage{titling}
\usepackage{titlesec}
\usepackage{amsmath}

\titleformat{\section}{\large\bfseries}{\thesection}{0.8em}{}
\titleformat{\subsection}{\normalsize\bfseries}{\thesubsection}{0.6em}{}

\title{\textbf{Argos Progress Report}\\
\large A multi-agent trading desk on four US stocks in 2024}
\author{Fuse-AI Fellowship submission}
\date{October 2026}

\begin{document}
\maketitle

\section{What we built, in one paragraph}

We built a trading research desk made of eight small language-model workers.
Three of them read the news, the earnings, and the price chart. Two of them
argue --- one trying to make the case for buying, one for selling. A judge
reads the argument without knowing which side wrote which, and picks the
stronger case. A risk manager checks the whole thing and can veto the trade.
A portfolio manager writes the final decision. The whole chain runs once per
week per stock, and we ran it across all 52 weeks of 2024, on four stocks:
Nvidia, Tesla, Oracle, and Coca-Cola. That is 212 full decisions in total.

Alongside the agents, a small neural network called Kronos reads the recent
price bars and predicts where the price is likely to go over the next 10
trading days. The agents see Kronos's forecast as one of their inputs.

Every decision ends up as a saved file on disk: all the inputs, every word
the agents said, how confident each one was, how many tokens the AI used,
and the final buy/sell/hold call. That saved file is the raw material for
everything you are about to read.

\section{The headline numbers}

We started each stock with a paper account of \$1{,}000. Four stocks, so
\$4{,}000 total. Here is what we ended 2024 with.

\begin{center}
\begin{tabular}{lrrrr}
\toprule
Stock & Starting cash & Ending equity & Profit (\$) & Return \\
\midrule
NVDA & \$1{,}000.00 & \$1{,}018.38 & \textcolor{black!70!green}{+\$18.38} & +1.84\% \\
TSLA & \$1{,}000.00 & \$1{,}000.00 & \$0.00 & 0.00\% \\
ORCL & \$1{,}000.00 & \$1{,}000.23 & +\$0.23 & +0.02\% \\
KO   & \$1{,}000.00 & \$1{,}002.28 & +\$2.28 & +0.23\% \\
\midrule
\textbf{Total} & \textbf{\$4{,}000.00} & \textbf{\$4{,}020.89} & \textcolor{black!70!green}{\textbf{+\$20.89}} & \textbf{+0.52\%} \\
\bottomrule
\end{tabular}
\end{center}

The desk made money, but not a lot. A reasonable question is: why so little?
Three honest answers.

\textbf{First}, the position sizing is deliberately small. The risk manager
caps each trade at 20\% of the account's cash. On a \$1{,}000 book that is
\$200 per trade. Even on NVDA, which rallied 171\% in 2024, the pipeline
only had \$200 at work at a time.

\textbf{Second}, the simulator was using fees originally written for the
Nepal stock market, which are higher than NASDAQ's actual fees. It took
0.40\% commission plus 7.5\% capital-gains tax out of every winning trade.
Fixing this later would add a few cents back on each stock.

\textbf{Third}, the system is cautious. It said \textit{hold} far more
often than it said \textit{buy}. On TSLA it never actually opened a trade,
because the two advocates were never confident enough on the same side.
That is why TSLA's ending equity is exactly \$1{,}000 --- the cash just
sat there the whole year.

The right way to read this: we were testing \textit{whether the system
calls direction correctly}, not whether it maximises profit. That is a
later knob to turn, once we trust the directional calls.

\section{Did the cache give some stocks an unfair edge?}

Fair question --- NVDA and TSLA were re-run many times (daily cadence, then
weekly, then interrupted and resumed), while ORCL got some re-runs and KO
was fully fresh with zero cache hits. If cached answers were somehow better
than fresh ones, the comparison would be unfair.

\textbf{They weren't.} Here is why:

The cache is keyed on \texttt{(agent name, model, exact prompt, decision
date)}. When a cache hit happens, the system replays the exact AI response
it got the first time. It does not peek at the future. It does not know
what the stock did later. It simply avoids re-paying for a response it
already has.

What \textit{could} be different between a cached ticker and a fresh one
is response variability. The AI runs at temperature 0.3, so each fresh
call picks slightly different words. A cached call returns one frozen
answer every time. So cached-tickers are more reproducible across runs,
and fresh ones have some sampling noise. That is not an advantage --- it
is just stability.

The clearest proof is KO itself. KO had zero cache hits. Every single one
of its 53 decisions was computed fresh. KO's final decision accuracy at
one week is 81.8\% --- essentially identical to NVDA's 82.4\%. If the
cache were secretly inflating NVDA, KO's fresh number would be lower.
It is not.

\section{What each metric means, in plain words}

Six numbers run the show. Here is what each one is actually measuring.

\subsection*{Directional accuracy at $h=5$}

If the system says \textit{buy} on a given Tuesday, and the stock's closing
price five trading days later is \textit{higher} than today's close, the
call counted as right. If it says \textit{sell} and the price was lower,
also right. The \textit{hold} calls are set aside --- hold is not a
directional prediction, it is a refusal to make one.

Accuracy is simply the fraction correct out of all directional calls.
\textbf{82\%} on NVDA means that of all the buy/sell calls the system
made for NVDA, 82 out of 100 pointed the right way one week later.
Coin-flip would be 50\%. Anything above 60\% is doing real work.

We measure the same thing at one day ($h=1$), one week ($h=5$), and one
month ($h=20$). On this run, the accuracy got \textit{better} as the
horizon got longer. That tells us the system is picking up slow,
multi-week signal, not just daily noise.

\subsection*{Sharpe ratio}

Imagine two traders who both made \$100 last year. One did it with a
smooth equity curve that climbed steadily. The other did it with wild
up-and-down swings that happened to end at the same spot. The second one
took on more risk for the same reward. Sharpe punishes the second one.

Formally, Sharpe is the average return divided by the standard deviation
of returns, annualised. In practice:

\begin{itemize}
\item Sharpe below 0: the strategy lost money or barely matched the swings.
\item Sharpe around 1: respectable --- the return paid for the swings.
\item Sharpe around 2: good --- real edge after accounting for volatility.
\item Sharpe above 3: unusual, often too good to be true.
\end{itemize}

NVDA's Sharpe of \textbf{1.44} says the return we got was worth the swings
we endured. Oracle's \textbf{1.10} is similar. KO's \textbf{0.68} is
modest --- Coca-Cola moves so slowly that even small swings dominate the
tiny profit.

\subsection*{Sortino ratio}

Sortino is Sharpe's smarter cousin. Sharpe treats all swings as bad ---
even the ones where you were winning. Sortino only penalises the
\textit{downward} swings. The reasoning: an investor doesn't care if their
account suddenly jumps up. They only care about the suddenly-down moves.

NVDA's Sortino of \textbf{2.16} and Oracle's \textbf{2.76} both say the
same thing: when we lost, we didn't lose much. The downside was tightly
controlled.

\subsection*{Max drawdown}

If you had been watching the paper account's balance tick by tick, this
is the single biggest percentage drop from a peak before it recovered.
It is the worst-case pain the account went through, not the worst-case
loss on a single trade.

NVDA's max drawdown was \textbf{$-$0.73\%}, meaning the account at its
lowest point was 0.73\% below its highest point. On a \$1{,}000 book
that's a worst-of-year drop of \$7.30. The pipeline hugs the risk floor.

\subsection*{Calmar ratio}

Return divided by the max drawdown. It answers: ``for every percent of
pain I took, how much return did I get?''

NVDA's Calmar of \textbf{2.49} says we made \$2.49 of return for every
\$1 of drawdown pain. Anything above 1 is reasonable; above 3 is strong.

\subsection*{Confidence gap}

This one is the research-critical number, and it needs the most explanation.

Both the Bull and the Bear advocate output a confidence between 0 and 1
along with their argument. If Bull averages 0.75 confidence across the
year and Bear averages 0.40, the \textit{confidence gap} is 0.35. On the
previous version of this project (built on the Nepal stock market), Bull
and Bear were running on \textit{different} AI models, and the gap was a
constant 0.30+ in every market. That gap wasn't measuring the market ---
it was measuring the fact that one model was more verbose than the other.

For this NASDAQ run we rewrote the router: both advocates share
\texttt{gpt-4o}, same temperature, same fallback chain. There is a
test that fails if the two configurations ever drift apart. If the gap
still shows up, it has to come from the evidence, not the models.

Here is what we saw across four very different 2024 markets:

\begin{center}
\begin{tabular}{llccc}
\toprule
Stock & 2024 character & Bull mean conf & Bear mean conf & Gap \\
\midrule
NVDA & clear bull year     & 0.679 & 0.419 & 0.260 \\
TSLA & bearish/chop        & 0.452 & 0.562 & 0.110 \\
ORCL & mixed               & 0.475 & 0.528 & 0.054 \\
KO   & flat (+5\%)         & 0.480 & 0.488 & 0.009 \\
\bottomrule
\end{tabular}
\end{center}

The gap collapses from 0.260 to 0.009 as the market goes from
one-sided-bullish to flat-and-boring. On TSLA, the \textit{Bear} is more
confident than the Bull. These are all the same two models, same prompts,
same temperature. The only thing that changed is the stock. That is proof
the debate reflects the market, not the AI.

\subsection*{Token usage and cost}

Every time an agent speaks, the AI reads a prompt and writes a response.
Each word is roughly a token (actually a sub-word fragment, but close
enough). We tracked exactly how many tokens went in and out at each step.

Across all 212 decisions:

\begin{itemize}
\item \textbf{5{,}601{,}846 tokens} total
\item \textbf{1{,}461 LLM calls} total
\item Average \textbf{26{,}400 tokens per decision}
\item Average \textbf{6.9 LLM calls per decision}
\end{itemize}

At gpt-4o pricing (roughly \$2.50 per million input tokens, \$10 per
million output), the whole study cost roughly \$6--8 in API fees.

The \textit{profit-per-1k-tokens} metric asks: for every 1{,}000 tokens
of AI thinking, how much paper P\&L did we generate? On a \$1{,}000 book
this looks tiny --- NVDA earns about \$0.013 of P\&L per 1k tokens. But
P\&L scales with book size while token spend doesn't. On a \$100{,}000
book the same pipeline would earn \$1.30 per 1k tokens, versus a token
cost of roughly \$0.0075. That is a 170$\times$ return on compute. The
economics only work on larger books.

\section{Stock-by-stock story}

\subsection*{NVDA --- the AI darling}

Nvidia had its year. Up 171\% in 2024, essentially straight line, driven
by data-centre AI demand. Our pipeline rode it with 82\% directional
accuracy at the one-week horizon, Sharpe 1.44, Sortino 2.16. The Bull
advocate was right 62\% of the time; the Bear 38\%. The confidence gap
was the widest of the four at 0.260 --- Bull was reading the same tape
Bear was, and judging the evidence as much stronger.

Where the pipeline left money on the table: it only opened one trade,
sized small, held to year-end. A more aggressive sizer would have put
\$500 or \$800 per trade and reloaded on dips.

\subsection*{TSLA --- the chop that trapped nobody}

Tesla's 2024 was a mess: down most of the year, then a late rally. Across
53 weekly decisions, the system only produced 4 directional calls. The
Portfolio Manager folded the chain into \textit{hold} 49 times. The paper
account sat at \$1{,}000 all year.

That looks like inactivity, but it is the right behaviour. On TSLA the
Bull and Bear were nearly tied in confidence (0.452 vs 0.562, Bear
slightly ahead), the technical readings were conflicting, and the
fundamentals wavered. When the ensemble refuses to commit, it holds.
Zero profit, zero drawdown. The Bear advocate was right 57\% of the time
on its calls --- more than the Bull.

\subsection*{ORCL --- the quiet AI play}

Oracle rallied on its cloud-and-AI capex guidance in 2024, but less
cleanly than NVDA. Our pipeline hit 50\% directional accuracy on 6
directional calls --- a small sample. Sharpe 1.10, Sortino 2.76. The
confidence gap was small (0.054) and leaned slightly bearish, which
matches the choppiness.

\subsection*{KO --- the sleepy legacy name}

Coca-Cola was the control: zero AI exposure, +5\% for the year, nothing
dramatic. If the pipeline only worked on AI-boosted names it would fail
here. It didn't. KO's final decision accuracy at one week was 81.8\%,
essentially matching NVDA. The confidence gap was 0.009 --- the
advocates were as close to perfectly tied as you can be.

This matters because KO was computed entirely fresh (no cache hits).
Every one of its 53 decisions was a real call to the AI. The 82\%
accuracy came from live reasoning, not replayed cache.

\section{What we did not do (yet)}

\begin{itemize}
\item \textbf{Fix the fee model for US stocks.} Still using Nepal retail
fees. Small change, meaningful on tight books.
\item \textbf{Scale position sizing.} Currently capped at 20\% of cash
per trade. A proper research setup would compare three sizing policies
side by side.
\item \textbf{Multi-round debate.} Each decision runs one round of Bull
vs Bear. More rounds, where each side can rebut the other's specific
claims, is a known upgrade --- not tested here.
\item \textbf{A benchmark portfolio.} We should compare the pipeline
against the obvious baselines: buy-and-hold, 60/40, equal-weight the
four tickers. All are easy to add to the eval script.
\item \textbf{Walk-forward instead of in-sample.} 2024 was both the
training window for the models and the backtest window. A proper
out-of-sample test would run on 2025 data.
\end{itemize}

\section{The honest caveat}

The AI models (OpenAI's gpt-4o and gpt-4o-mini) were trained on text that
includes 2024. When we ask the model to reason about NVDA ``as of
2024-01-02,'' it may already know what happened to NVDA later in 2024.
We tell it in the prompt to reason strictly as of the given date, but
there is no mechanical way to force an AI to forget what it knows.

This is a real limitation. The directional accuracy we see is probably
inflated by some amount of implicit knowledge. The right fix is to re-run
on 2026 data, where the models have no training signal yet. That is a
natural next step.

What is \textit{not} inflated by this is the symmetric-debate finding.
Both Bull and Bear have access to the same model with the same knowledge.
If Bear is more right than Bull on TSLA, that is because the evidence on
TSLA pointed down, not because one model knew more than the other.

\section{Where this leaves us}

The pipeline makes directional calls that are right about 80\% of the
time on three of four stocks, holds through contested markets, controls
drawdown tightly, and costs about \$6--8 of API fees to run across a
year of weekly decisions on four stocks. The debate is now
provably driven by evidence rather than model bias. The research
scaffolding --- journals, token tracking, dense NAV, Sharpe/Sortino/
Calmar/max-drawdown, and the eval harness --- is in place and can be
pointed at a longer window, more stocks, or a different market with
one config change.

The next real experiment is a 2026 out-of-sample run on five or six
names, with the fee model corrected and position sizing scaled up. The
code is ready.

\end{document}
